\section{Related Work}
\label{sec:related_work}

Incremental 2D aerial mosaicking from unmanned aerial vehicle (UAV) video requires addressing sequential registration, computational resource limits, failure detection, and error mitigation. We organize prior literature into four interrelated domains and establish the scope of our comparative evaluation.

\subsection{Incremental Aerial Mosaicking and Monocular SLAM}
Early and contemporary UAV video mosaicking systems rely predominantly on sequential homography chaining. Hwang et al.~\cite{hwang2026} developed a CPU-based real-time 2D orthomosaic framework using incremental homography estimation with classical feature detectors (SIFT~\cite{lowe2004}, ORB~\cite{rublee2011}) and RANSAC~\cite{fischler1981}, relying on inlier count thresholds to accept inter-frame transformations. Li et al.~\cite{li2023} proposed a real-time incremental video mosaic pipeline for UAV remote sensing incorporating weighted reprojection error optimization, overlap-based keyframe selection, and historical keyframe rematching to mitigate tracking drift. Chen et al.~\cite{chen2022} addressed perspective distortion across varying flight elevations by combining dynamic programming seamline estimation with shape-preserving half-projective warps. To overcome the geometric vulnerability of pure pairwise homography chains, Bu et al.~\cite{map2dfusion} introduced Map2DFusion, coupling direct monocular visual SLAM tracking with keyframe graph optimization and local bundle adjustment. In contrast to full visual SLAM frameworks that maintain 3D map points and solve bundle adjustment, our investigation focuses on lightweight 2D incremental homography pipelines operating under sequential resource limits without global graph optimization.

\subsection{Learned Feature Matching and Efficiency Constraints}
Deep learning has substantially advanced feature matching under challenging illumination and viewpoint changes. DeTone et al.~\cite{detone2018} introduced SuperPoint, a self-supervised framework for joint interest point detection and description. Lindenberger et al.~\cite{lindenberger2023} presented LightGlue, a graph neural network matcher that adaptively prunes unmatchable points and enables early stopping. For detector-free matching, Sun et al.~\cite{sun2021} developed LoFTR using Transformers with coarse-to-fine attention, while Wang et al.~\cite{wang2024} proposed Efficient LoFTR to reduce computational complexity through an aggregated attention mechanism with adaptive token selection and a two-stage correlation layer. For edge deployment, Potje et al.~\cite{xfeat2024} developed XFeat, an architecture tailored for lightweight matching on resource-constrained platforms. Kim and Kim~\cite{kim2025} recently combined classical SIFT extraction with LightGlue to handle large rotational displacements in UAV scenes.

While learned matchers demonstrate high pairwise registration capability, their sequential behavior in continuous UAV video chains exhibits practical trade-offs. In long flight sequences without global bundle adjustment, successive pairwise learned transformations can still accumulate residual subpixel errors across nominal flight segments. This motivates evaluating selective escalation strategies rather than continuous unconstrained invocation.

\subsection{Registration Failure Detection and Uncertainty Estimation}
Zaffar et al.~\cite{zaffar2024} compare retrieval-based uncertainty, learned aleatoric uncertainty, and geometric verification for visual place recognition, and introduce a baseline that uses reference-image poses. Their target is confidence in image retrieval, not accumulated aerial-mosaic drift; this distinction motivates evaluating indicators against an explicitly defined registration-error proxy. Barath et al.~\cite{barath2023} provide reference homographies for benchmarking robust estimators and correspondence filtering methods. Ivashechkin et al.~\cite{ivashechkin2021} introduce VSAC, which uses the concept of independent inliers and adaptive sequential hypothesis verification to improve robust estimation. These works distinguish confidence assessment and robust fitting from the separate problem of managing an accepted transformation over a long chain.

Internal solver outputs are useful acceptance signals~\cite{hwang2026,li2023}, but a small residual on selected correspondences does not by itself establish globally consistent registration. Our study examines whether post-registration appearance consistency and geometric indicators add descriptive evidence beyond point counts; it does not assume that photometric agreement proves geometric correctness.

\subsection{Sequential Error Propagation and Recovery Policies}
When registration degradation occurs, an incremental mapping pipeline requires a recovery strategy to handle the corrupted step. The simplest policy, frame dropping, discards unreliable frames~\cite{hwang2026}. However, discarding frames dilates the temporal baseline between incoming frames and the active mosaic anchor, increasing perspective disparity and potentially triggering subsequent drops. Li et al.~\cite{li2023} addressed this by searching past keyframes for rematching candidates, though successful reconnection depends on spatial trajectory overlap and search window constraints. Kharismawati and Kazic~\cite{kharismawati2025} developed DroneZaic for agricultural UAV videos, using motion clustering to partition long flights into independent submosaics.

Submosaic splitting preserves the entire raw frame inventory by opening new reference frames during severe registration disruptions, though this trades single-canvas metric continuity for localized frame coverage. Conversely, frame dropping preserves continuity within surviving segments but risks discarding substantial portions of the flight when degradation persists. In this work, we compare these recovery policies under unified risk-triggering criteria.
